\paragraph{Conjugate-coordinate class: Lorentzian moments in the fibre-bounded cone window, and the vector, ghost and spin terms of the invariant sector (conditional; bounded no-go).}
Keep the setting and notation of the two preceding paragraphs (Poisson limit, leaf \(\mathbb{CP}^{1,2}\), \(T^a=M^{a4}\), \(\Box\), \(H_V\) with \(\mu^2=-3\), \(C_{42}\), \(C_{41}\), \(C_{31}\), \(\sigma(\cdot)\), \(S^2=2C_{41}-C_{42}\) with symbol \(\beta q_f\), the base symbol \(\alpha q_b\), \(a=\sinh\tau=y^4/\ell\), and the function \(Q(a)\) of the exact invariant sector). At a point of the leaf take coordinates \((u,\zeta)\in\mathbb R^4\times\mathbb R^2\) on the cotangent space with \(|u|^2=\alpha q_b\) and \(|\zeta|^2=\beta q_f\), so that \(-\sigma(C_{42})=|u|^2-|\zeta|^2\), and write \(\mathsf K\) for the form \(\sigma(\Box)\) (rank four, signature \((3,1)\)) in these coordinates, with blocks \(\mathsf K_{uu}\), \(\mathsf K_{u\zeta}\), \(\mathsf K_{\zeta\zeta}\).

\emph{Gaussian windows.} Let \(\lambda_0=-c\) and \(\lambda_1,\lambda_2,\lambda_3\) be positive and \(\mathcal M=\mathrm i\int_0^\infty\mathrm ds\prod_{k=0}^3(1+\mathrm is\lambda_k)^{-1/2}\), the Gaussian average of \((\sigma-\mathrm i0)^{-1}\). Rotating \(s=-\mathrm iu\) (the fourth quadrant contains no singularity) gives \(\mathcal M=\int_0^\infty\mathrm du\,(1-cu-\mathrm i0)^{-1/2}\prod_{k\ge1}(1+\lambda_ku)^{-1/2}\), whose first factor is real for \(u<1/c\) and equals \(\mathrm i(cu-1)^{-1/2}\) for \(u>1/c\). With \(u=(1\mp v^2)/c\),
\[
  \operatorname{Re}\mathcal M=\frac2c\int_0^1\!\mathrm dv\prod_{k\ge1}\Bigl(1+\frac{\lambda_k(1-v^2)}c\Bigr)^{-1/2}>0,\qquad
  \operatorname{Im}\mathcal M=\frac2c\int_0^\infty\!\mathrm dv\prod_{k\ge1}\Bigl(1+\frac{\lambda_k(1+v^2)}c\Bigr)^{-1/2}>0 ,
\]
and \(\operatorname{Re}\mathcal M=2/(\lambda+c)\) for \(\lambda_k=\lambda\). The pencil of \(\sigma(\Box)\) with any positive definite window form has one negative, three positive and two zero generalised eigenvalues, so the principal-value mass moment of every Gaussian window of the class is positive; this proves, for Gaussian windows, the positivity that the preceding paragraph found by sampling.

\emph{Geometry of the cone window.} At \(60\) random orbit points (agreement \(10^{-9}\) or better, not proved): the eigenvalues of \(\mathsf K_{uu}\) are \(-\tfrac12\) and \(\tfrac12\sinh^2\tau\) (three times); the Schur complement of \(\mathsf K_{uu}\) in \(\mathsf K\) vanishes, so \(\sigma=Q_L(u+w\zeta)\) with \(Q_L=\mathsf K_{uu}\) and \(w=\mathsf K_{uu}^{-1}\mathsf K_{u\zeta}\); \(w\) is conformal, \(|w\zeta|=\coth\tau\,|\zeta|\), and its image is orthogonal to the time eigenvector of \(Q_L\). In the variable \(\boldsymbol x=u+w\zeta\) the removed region \(|u|<|\zeta|\) is therefore the ball of radius \(r=|\zeta|\) about the spacelike point \(p=w\zeta\), at distance \(r\coth\tau=r\sqrt{1+c/\lambda}\) from the apex \(\boldsymbol x=0\) with \(c=\tfrac12\), \(\lambda=\tfrac12\sinh^2\tau\): the cone \(|u|=|\zeta|\) touches the null cone of \(Q_L\).

\emph{Mass moment in the fibre-bounded cone window.} Let \(\chi\ge0\) be bounded and integrable on \([0,\infty)\) and take the window function \(\chi(|u|^2-|\zeta|^2)\) on \(|\zeta|^2\le Q_f\) and \(|u|\ge|\zeta|\). For a ball of radius \(R\) about \(p\) with \(|p|<R\), polar coordinates about the apex give \(\operatorname{PV}\!\int_{|\boldsymbol x-p|<R}\mathrm d^4\boldsymbol x\,Q_L^{-1}=\tfrac12(R^2-|p|^2)I_0+p^{\mathsf T}\mathsf Np\), where \(I_0=\operatorname{PV}\oint\mathrm d\Omega\,Q_L^{-1}=8\pi^2/\cosh^2\tau\) and \(\mathsf N=\operatorname{PV}\oint\mathrm d\Omega\,nn^{\mathsf T}Q_L^{-1}\), because the remaining term of the exit radius is odd in the direction \(n\) and \(Q_L\) is even (checked against a direct quadrature to \(10^{-4}\)). For spacelike \(p\) with \(|p|\ge R\), quadrature indicates that the same expression continues to hold up to the tangency distance \(R\coth\tau\) (agreement to \(10^{-3}\) in the sampled cases) and departs from it beyond; this is a numerical observation and is not proved here. Under this observation every shell \(r\le|u|\le\sqrt{r^2+\Lambda^2}\) has principal-value moment \(\tfrac12\Lambda^2I_0\), independent of \(p\), and so
\[
  \operatorname{Re}\,\mathcal N_W=\tfrac12\,\varrho\,\pi Q_f\,I_0\int_0^\infty\chi(\lambda)\,\mathrm d\lambda>0 ,\qquad
  \mathcal N_W=\int_W\varrho\,\mathrm d^4u\,\mathrm d^2\zeta\;\chi\,(\sigma-\mathrm i0)^{-1},
\]
with \(\varrho>0\) the Jacobian of the coordinates \((u,\zeta)\); the sign of the real part is conditional on the numerical observation, and the real part depends on the cutoff function only through \(\int\chi\) and on the fibre bound only through the area \(\pi Q_f\). The imaginary part is \(\pi\int_W\varrho\,\chi\,\delta(\sigma)\ge0\). Finiteness of \(\mathcal N_W\) is likewise rested on the shell formula: the apex singularity is of order \(|\boldsymbol x|^{-2}\) in four dimensions and integrable, but the window function has a jump at the cone edge when \(\chi(0)\ne0\), and the edge touches the null cone of \(Q_L\) at one point, so finiteness there is not shown independently.

\emph{Zeeman field and electric growth.} The first-order term of \(H_V\) is exactly the Zeeman matrix of the field \(F_{ab}(\xi)=\xi(\mathcal J_{ab})\), linear in the covector \(\xi\); let \(\eta F\) have eigenvalues \(\pm f_e\) (electric pair, factor \(\mathrm e^{\pm2sf_e}\) in \(\mathrm e^{-\mathrm isH_V}\)) and \(\pm\mathrm if_m\) (magnetic pair, phases), so \(f_e^2-f_m^2=-F_{ab}F^{ab}/2=-\sigma(C_{31})(\xi)\). At \(59\) of the \(60\) points (one point is numerically singular in the whitening, a fibre Gram eigenvalue vanishing in double precision), to \(10^{-9}\) except for the maximum over the base plane, which is obtained from \(3000\) sampled directions to \(6\cdot10^{-5}\): on the base plane \(f_m=0\) and \(\max f_e=|u|\cosh\tau/\sqrt2\); on the vertical plane of the pencil \(f_e=|\zeta|\sinh\tau/\sqrt2\), \(f_m=|\zeta|/\sqrt2\) in every direction; on the kernel plane of \(\mathsf K\) (the centre variable of the nilpotent model) \(f_e^2=|\zeta|^2/(2(1+2\sinh^2\tau))\) and \(f_m^2=|\zeta|^2\sinh^2\tau/(2(1+2\sinh^2\tau))\), so that on this plane \(F_{ab}F^{ab}/2\) is negative for \(|y^4|<\ell\), positive for \(|y^4|>\ell\) and changes sign exactly at \(|y^4|=\ell\); this exact threshold replaces the sample-dependent value of the preceding paragraph on that plane (on the vertical plane the signs are opposite). The fibre-driven electric rate \(2f_e\) is \(2s_f\sinh\tau\) with \(\beta q_f=2s_f^2\le2s_{\max}^2\), whereas the base-driven rate grows linearly with \(|u|\); the electric rate is thus bounded by a quantity linear in the base cutoff, whose fibre part is bounded independently of the base cutoff. The logarithm of the determinant of the Zeeman block, \(\ln(\sigma^2-4f_m^2)+\ln(\sigma^2+4f_e^2)\), has only integrable logarithmic singularities, at \(\sigma=\pm2f_m\), at the leading symbol; the growth is a property of the Schwinger representation and is not removed by the fibre bound.

\emph{Vector and ghost sectors of the invariant sector.} The ghost operator is \(-\Box\), the same as the scalar of the preceding paragraph. On fibre-constant vector fields \(A_c(y)\) the cross term \(\eta_{ac}\eta_{bd}\mathcal J^{ab}\Sigma^{cd}\) of \(C_G\) equals \(-2s\mathcal J_{ab}A^b\) exactly (symbolic check), so \(H_V=-\Box-2s\mathcal J_{ab}A^b+\mu^2\) is \(C_G-C_{41}-C_\Sigma+\mu^2\) there. Write the diagonal of \(\mathrm e^{z\Delta_{H^4}+w\Delta_{H^3}}\) as \(D=D_0(1+D_1+\dots)\), with \(D_1\) depending on \(a\) and homogeneous of degree one in \((z,w)=(t-\mathrm is,\mathrm is)\); \(\operatorname{tr}(E+R/6)\) is the first Seeley--DeWitt density of \(\mathrm e^{t\Delta_{H^4}-\mathrm isH_V}\) in the convention that the exponent is \(\Delta+E\), and for one scalar component it equals \(D_1\). For the four components, with the \(4\times4\) matrix connection,
\[
  \operatorname{tr}\bigl(E+\tfrac16R\bigr)=4D_1+\frac{6s^2(1+a^2)}{t+\mathrm isa^2}-4\mathrm is\mu^2 ,
\]
the spin term being \(2s^2\operatorname{tr}[G\,\sigma(\Delta_{H^3})]=-2s^2\operatorname{tr}[G\,\sigma(C_{31})]\), with \(G\) the inverse of the symbol of the exponent, that is the Zeeman square \(f_e^2-f_m^2\) traced against \(G\). The route reproduces the exact scalar coefficient \(D_1\) to \(120\) digits at three points, and the vector and Dirac traces were verified to the same accuracy; the vector kernel itself was not computed. With \(x=s/t\) and the subtraction at \(x=1\) of the preceding paragraph, the weight-two functions are \(Q_\mu=-\mu^2j(a)\) with \(j(a)=2(1+\mathrm i/a)/(1+a^2)\), and \(Q_{\mathrm{sp}}=-8/(1+a^2)-\mathrm i\,4(3a^2+1)/(a^3(1+a^2))\) for the spin term, so \(\operatorname{Re}[Q_{\mathrm{sp}}+4Q_{\mu}]=16/(1+a^2)\) for \(\mu^2=-3\). For \(n_{\mathrm{ex}}\) further scalars of mass \(\mu_i^2\) the bosonic content has
\[
  \operatorname{Re}Q_{\mathrm{bos}}=(2+n_{\mathrm{ex}})\operatorname{Re}Q+\frac{16-2n_{\mathrm{ex}}\mu_i^2}{1+a^2},
\]
plus the scale-dependent term proportional to \(2+n_{\mathrm{ex}}\) of the preceding paragraph. The first term is indefinite and the second has the sign of \(8-n_{\mathrm{ex}}\mu_i^2\); the scale-dependent term has the sign of a logarithm. On the full leaf the spin term in a Gaussian window has a large-\(s\) tail \(C\,\mathrm ds/s\) with \(C\) purely imaginary and proportional to \(\sigma(C_{31})\) on the kernel plane, that is to \((1-a^2)/(1+2a^2)\), so its logarithmic coefficient is of indefinite sign on the leaf, and its real part depends on the contour angle.

\emph{Supertrace.} The matrix part of the coefficient is traceless, so it enters only through \(\operatorname{tr}\Sigma^{ab}\Sigma^{cd}\), and the species-independent part \(R/6+E\) is multiplied by the number of components. For the Dirac representation \(\operatorname{tr}\Sigma_D\Sigma_D=\tfrac12\operatorname{tr}\Sigma_V\Sigma_V\) for all index pairs. Supply, as an input not contained in the class, the sixteen real fermion components of the ten-dimensional supersymmetric matrix model with the square of the Dirac operator equal to \(\Box+s\Sigma^{ab}\mathcal J_{ab}\) on each component, no bilinear mass, and \(n_{\mathrm{ex}}=6\), and weight \(-\tfrac12\) per real fermion component. Then the component count \(4+6-2-8\) vanishes, the weighted index \(1-\tfrac12\cdot4\cdot\tfrac12\) vanishes, the Zeeman character is \(4(\cos sf_m-\cosh sf_e)^2=s^4(f_e^2+f_m^2)^2+\dots\) (the order \(s^4\) is a sum of the squares of the two invariants and does not vanish for any non-zero flux in Lorentzian signature), and what remains of the weight-two coefficient in the invariant sector is \(-\bigl(4\mu^2+n_{\mathrm{ex}}\mu_i^2\bigr)j(a)\), the mass supertrace times the mass moment \(j(a)\), with \(\operatorname{Re}j=2/(1+a^2)>0\); no sign of the mass supertrace itself is claimed. The mode-count term and the spin term of the order \(F^2\) both cancel identically; the curvature and connection part of the weight-two term therefore does not by itself carry the Newton coefficient for this content; whether it arises at higher order of the expansion or through the dependence on \(A\) is not computed.

This is a statement for the Poisson limit, one loop, the leading symbol and the invariant sector, the four-matrix class with \(\mu^2=-3\), Gaussian windows and the cone window of the stated form, the sampled points and the supplied fermion content with no bilinear fermion mass and the same window for bosons and fermions. The continuation of the shell identity to spacelike \(p\) with \(|p|\ge R\) is a numerical observation, and the sign of the cone-window moment is conditional on it; the vector kernel and the second variation under \(A\) are not computed; fermions are not part of the class; the Jacobian \(\varrho\) and the Liouville normalisation are not computed as functions of position. It does not treat Moyal terms or a window built from dynamical matrices, it identifies no induced metric with the emergent metric of Chapter~\ref{ch:vierbein-emergence} and, since a Euclidean coefficient is not a Lorentzian one, leaves Assumption~\ref{ass:lorentzian} in place (Problem~E). This changes no ledger label; Problem~F remains open.
